\documentclass[10pt,a4paper]{amsart}
\usepackage{amsmath,amssymb,amsthm,graphicx}
\usepackage{verbatim}
\newenvironment{reference}{\comment}{\endcomment}
%\newenvironment{reference}{}{}
\newenvironment{slogan}{\comment}{\endcomment}
\newenvironment{history}{\comment}{\endcomment}

% For commutative diagrams we use Xy-pic
\usepackage[all]{xy}

% We use 2cell for 2-commutative diagrams.
\xyoption{2cell}
% Theorem environments.
%
\theoremstyle{plain}
\newtheorem{theorem}[subsection]{Theorem}
\newtheorem{proposition}[subsection]{Proposition}
\newtheorem{lemma}[subsection]{Lemma}

\theoremstyle{definition}
\newtheorem{definition}[subsection]{Definition}
\newtheorem{example}[subsection]{Example}
\newtheorem{exercise}[subsection]{Exercise}
\newtheorem{situation}[subsection]{Situation}

\theoremstyle{remark}
\newtheorem{remark}[subsection]{Remark}
\newtheorem{remarks}[subsection]{Remarks}

\numberwithin{equation}{subsection}

% Macros
%
\def\lim{\mathop{\mathrm{lim}}\nolimits}
\def\colim{\mathop{\mathrm{colim}}\nolimits}
\def\Spec{\mathop{\mathrm{Spec}}}
\def\Hom{\mathop{\mathrm{Hom}}\nolimits}
\def\Ext{\mathop{\mathrm{Ext}}\nolimits}
\def\SheafHom{\mathop{\mathcal{H}\!\mathit{om}}\nolimits}
\def\SheafExt{\mathop{\mathcal{E}\!\mathit{xt}}\nolimits}
\def\Sch{\mathit{Sch}}
\def\Mor{\mathop{\mathrm{Mor}}\nolimits}
\def\Ob{\mathop{\mathrm{Ob}}\nolimits}
\def\Sh{\mathop{\mathit{Sh}}\nolimits}
\def\NL{\mathop{N\!L}\nolimits}
\def\CH{\mathop{\mathrm{CH}}\nolimits}
\def\proetale{{pro\text{-}\acute{e}tale}}
\def\etale{{\acute{e}tale}}
\def\QCoh{\mathit{QCoh}}
\def\Ker{\mathop{\mathrm{Ker}}}
\def\Im{\mathop{\mathrm{Im}}}
\def\Coker{\mathop{\mathrm{Coker}}}
\def\Coim{\mathop{\mathrm{Coim}}}

% Boxtimes
%
\DeclareMathSymbol{\boxtimes}{\mathbin}{AMSa}{"02}

%
% Macros for moduli stacks/spaces
%
\def\QCohstack{\mathcal{QC}\!\mathit{oh}}
\def\Cohstack{\mathcal{C}\!\mathit{oh}}
\def\Spacesstack{\mathcal{S}\!\mathit{paces}}
\def\Quotfunctor{\mathrm{Quot}}
\def\Hilbfunctor{\mathrm{Hilb}}
\def\Curvesstack{\mathcal{C}\!\mathit{urves}}
\def\Polarizedstack{\mathcal{P}\!\mathit{olarized}}
\def\Complexesstack{\mathcal{C}\!\mathit{omplexes}}
% \Pic is the operator that assigns to X its picard group, usage \Pic(X)
% \Picardstack_{X/B} denotes the Picard stack of X over B
% \Picardfunctor_{X/B} denotes the Picard functor of X over B
\def\Pic{\mathop{\mathrm{Pic}}\nolimits}
\def\Picardstack{\mathcal{P}\!\mathit{ic}}
\def\Picardfunctor{\mathrm{Pic}}
\def\Deformationcategory{\mathcal{D}\!\mathit{ef}}
\usepackage{cleveref}
\crefname{originalremark}{original remark}{original remarks}
\Crefname{originalremark}{Original remark}{Original remarks}
\makeatletter
\newlabel{remark-QCoh-abelian}{{0GQL}{source}{}{sourcepointer-1}{}}
\newlabel{remark-QCoh-abelian@cref}{{[originalremark][0][]0GQL}{[1][0][]source}{}{}{}}
\makeatother
\crefname{originallemma}{original lemma}{original lemmas}
\Crefname{originallemma}{Original lemma}{Original lemmas}
\makeatletter
\newlabel{lemma-coherent-Noetherian}{{0GR5}{source}{}{sourcepointer-2}{}}
\newlabel{lemma-coherent-Noetherian@cref}{{[originallemma][0][]0GR5}{[1][0][]source}{}{}{}}
\makeatother
\crefname{originallemma}{original lemma}{original lemmas}
\Crefname{originallemma}{Original lemma}{Original lemmas}
\makeatletter
\newlabel{lemma-parasitic}{{0774}{source}{}{sourcepointer-3}{}}
\newlabel{lemma-parasitic@cref}{{[originallemma][0][]0774}{[1][0][]source}{}{}{}}
\makeatother
\crefname{originallemma}{original lemma}{original lemmas}
\Crefname{originallemma}{Original lemma}{Original lemmas}
\makeatletter
\newlabel{lemma-adjoint-kernel-parasitic}{{0779}{source}{}{sourcepointer-4}{}}
\newlabel{lemma-adjoint-kernel-parasitic@cref}{{[originallemma][0][]0779}{[1][0][]source}{}{}{}}
\makeatother
\title{Coherent modules on a locally Noetherian stack form an abelian category}
\author{The Stacks Project Authors}
\date{stacks-project:stacks-cohomology.tex:local-snapshot-20261012; GNU FDL 1.2 or later}
\begin{document}
\maketitle
\noindent\textit{Editorial note.} All delivered mathematical text is the author's own, copied verbatim from the pinned Stacks Project chapter file stacks-cohomology.tex; no statement or proof step was rewritten, reordered or removed.\par\medskip
\noindent\textit{Editorial note.} The counted span cites the lemmas recorded in assembly.external\_references with canonical tags as navigation pointers.\par\medskip
\noindent\textit{Allowed standard prerequisites (editorial).}\begin{itemize}
\item QCoh rules for kernels and cokernels. Cohomology of Algebraic Stacks, Remark stacks-cohomology-remark-QCoh-abelian. Prescription agreement.
\item Four coherent characterizations. Cohomology of Algebraic Stacks, Lemma stacks-cohomology-lemma-coherent-Noetherian. Coherence check on the presentation.
\item Adjunction defect is parasitic. Cohomology of Algebraic Stacks, Lemma stacks-cohomology-lemma-adjoint-kernel-parasitic (2). f\textasciicircum{}*Q(Ker) = f\textasciicircum{}*Ker.
\end{itemize}\medskip
\ifdefined\hypertarget\hypertarget{sourcepointer-1}{}\fi
\noindent\textit{Source reference.} Original remark 0GQL: sources/stacks-chow/tags.txt, source lines 20866--20866. Cross-chapter standard prerequisite cited by the author inside the counted span; it is used as a pointer and its own proof is not part of this item.\par\medskip
\ifdefined\hypertarget\hypertarget{sourcepointer-2}{}\fi
\noindent\textit{Source reference.} Original lemma 0GR5: sources/stacks-chow/tags.txt, source lines 20885--20885. Cross-chapter standard prerequisite cited by the author inside the counted span; it is used as a pointer and its own proof is not part of this item.\par\medskip
\ifdefined\hypertarget\hypertarget{sourcepointer-3}{}\fi
\noindent\textit{Source reference.} Original lemma 0774: sources/stacks-chow/tags.txt, source lines 9189--9189. Cross-chapter standard prerequisite cited by the author inside the counted span; it is used as a pointer and its own proof is not part of this item.\par\medskip
\ifdefined\hypertarget\hypertarget{sourcepointer-4}{}\fi
\noindent\textit{Source reference.} Original lemma 0779: sources/stacks-chow/tags.txt, source lines 9194--9194. Cross-chapter standard prerequisite cited by the author inside the counted span; it is used as a pointer and its own proof is not part of this item.\par\medskip
\begin{lemma}
\label{lemma-coherent-abelian-Noetherian}
Let $\mathcal{X}$ be a locally Noetherian algebraic stack.
The category of coherent $\mathcal{O}_\mathcal{X}$-modules is abelian.
If $\varphi : \mathcal{F} \to \mathcal{G}$ is a map
of coherent $\mathcal{O}_\mathcal{X}$-modules, then
\begin{enumerate}
\item the cokernel $\Coker(\varphi)$ computed in
$\textit{Mod}(\mathcal{O}_\mathcal{X})$ is a coherent
$\mathcal{O}_\mathcal{X}$-module,
\item the image $\Im(\varphi)$ computed in
$\textit{Mod}(\mathcal{O}_\mathcal{X})$ is a coherent
$\mathcal{O}_\mathcal{X}$-module, and
\item the kernel $\Ker(\varphi)$ computed in
$\textit{Mod}(\mathcal{O}_\mathcal{X})$
may not be coherent, but it is
in $\textit{LQCoh}^{fbc}(\mathcal{O}_\mathcal{X})$ and $Q(\Ker(\varphi))$
is coherent and is the kernel of $\varphi$ in
$\textit{Coh}(\mathcal{O}_\mathcal{X})$.
\end{enumerate}
The inclusion functor $\textit{Coh}(\mathcal{O}_\mathcal{X}) \to
\QCoh(\mathcal{O}_\mathcal{X})$ is exact.
\end{lemma}

\begin{proof}
The rules given for taking kernels, images, and cokernels in
$\textit{Coh}(\mathcal{O}_\mathcal{X})$ agree with the prescription
for quasi-coherent modules in Remark \ref{remark-QCoh-abelian}.
Hence the lemma will follow if we can show that the quasi-coherent modules
$\Coker(\varphi)$, $\Im(\varphi)$, and $Q(\Ker(\varphi))$ are coherent.
By Lemma \ref{lemma-coherent-Noetherian} it suffices to prove this after
restricting to $U_\etale$ for some surjective smooth morphism
$f : U \to \mathcal{X}$. The functor
$\mathcal{F} \mapsto f^*\mathcal{F}|_{U_\etale}$
is exact. Hence $f^*\Coker(\varphi)$ and $f^*\Im(\varphi)$ are
the cokernel and image of a map between coherent $\mathcal{O}_U$-modules
hence coherent as desired. The functor
$\mathcal{F} \mapsto f^*\mathcal{F}|_{U_\etale}$
kills parasitic modules by Lemma \ref{lemma-parasitic}.
Hence $f^*Q(\Ker(\varphi))|_{U_\etale} = f^*\Ker(\varphi)|_{U_\etale}$
by part (2) of Lemma \ref{lemma-adjoint-kernel-parasitic}.
Thus we conclude that $Q(\Ker(\varphi))$ is coherent in the same way.
\end{proof}

\end{document}
